%! Choosing and Comparing Models in Context — Unit 2, Lesson 2.7 — Lesson Plan
\documentclass[10pt]{article}
\usepackage{atda-article}
\usepackage{atda-boxes}
\usepackage{graphicx}   % -article does NOT load graphicx; needed for thumbnails/screenshots
\graphicspath{{images/}}
\newcommand{\TallMath}[1]{$\displaystyle #1\rule[-1.4em]{0pt}{3.2em}$}
\newcommand{\UnitNumberName}{Unit 2: Functions, Transformations, and Graph Analysis \quad}
\newcommand{\LessonNumberName}{Lesson 2.7: Choosing and Comparing Models in Context}

\begin{document}
\begin{center}
    {\Large\bfseries \CourseName \\
    {\small\bfseries \UnitNumberName \LessonNumberName}}
\end{center}

\begin{tcolorbox}[colback=mist, colframe=royal, boxrule=0.9pt, arc=2mm,
    left=2mm, right=2mm, top=1.5mm, bottom=1.5mm]
    \textbf{Primary Objective:} I can decide whether a \textbf{linear}, \textbf{quadratic}, or
    \textbf{exponential} function best models a table, a graph, or a story, and name the test that
    decided it --- first differences, second differences, or ratios from a table, and a straight edge,
    a turning point, or a flattening end from a graph. I can compare the three families across every
    representation I have, use the model I chose to answer a question in context both graphically and
    algebraically, and say where the model's answers stop being trustworthy.
\end{tcolorbox}

\renewcommand{\arraystretch}{1.6}
\begin{skillbox}[Priority Ideas \& Skills]{goldbox}
\begin{tabularx}{\linewidth}{@{}X|X@{}}
\begin{minipage}[t]{\linewidth}\vspace{0pt}
\textbf{Standards covered:}
\begin{itemize}[leftmargin=1.1em, itemsep=2pt, topsep=2pt]
  \item \texttt{AFDA.AF.1c} --- determine and analyze whether a linear, quadratic, or exponential
        function \textbf{best models} a given representation, including in context. This is the whole
        lesson.
  \item \texttt{AFDA.AF.1e} --- use a graphical or algebraic representation of a function to
        \textbf{solve problems within a context}, both ways, as appropriate.
  \item \texttt{AFDA.AF.1g} --- \textbf{compare and contrast} the three families using multiple
        representations: graphs, tables, equations, and verbal descriptions.
  \item \texttt{AFDA.AF.2h} --- \emph{relate} the characteristics of the graphs across families,
        including in context. (2.6 \emph{described} them; today they get compared.)
\end{itemize}
\textbf{Skills students leave with:}
\begin{itemize}[leftmargin=1.1em, itemsep=2pt, topsep=2pt]
  \item Run the table tests \textbf{in order} --- first differences, then second differences, then
        ratios --- on \textbf{equal} input steps.
  \item Choose a family off a \textbf{graph} by naming a feature, not a vibe.
  \item Choose a family off a \textbf{story} by hearing ``the same amount added'' versus ``the same
        percent.''
  \item Answer a contextual question from the chosen model, graphically \emph{and} algebraically.
  \item Say where the model stops being believable, and why.
\end{itemize}
\end{minipage}
&
\begin{minipage}[t]{\linewidth}\vspace{0pt}
\textbf{Key Understandings (the \emph{why}):}
\begin{itemize}[leftmargin=1.1em, itemsep=2pt, topsep=2pt]
  \item \textbf{The family is chosen by how the outputs \emph{change}, not by what the graph looks
        like.} That is the sentence the whole lesson runs on, and it is what makes the tests
        necessary rather than optional.
  \item \textbf{``It curves upward'' rules out linear and nothing else.} Quadratic and exponential
        both have growing first differences. This single confusion is the hook, the exit ticket, and
        two homework items.
  \item \textbf{A model is a choice, not a fact.} The guidance's own phrase is
        \emph{mathematical model} (\texttt{AFDA.DA.1d}); a different choice can fit the same data.
  \item \textbf{Two points never decide a family, and three points never rule out a quadratic.}
        Agreement on the data you have is not evidence.
  \item \textbf{Every model has a window of inputs where it works.} Reading past it is where a wrong
        family choice --- or a right one --- produces nonsense.
\end{itemize}
\end{minipage}
\end{tabularx}
\end{skillbox}

\begin{skillbox}[{Vocabulary, Concepts, \& Theorems}]{greenbox}
\renewcommand{\arraystretch}{1.35}
\begin{tabularx}{\linewidth}{@{}p{4.1cm}X@{}}
\textbf{Term} & \textbf{Meaning students record} \\
\hline
Mathematical model & A function chosen to describe a real situation. It is a \emph{choice}, not a
    fact --- and it can be wrong. \\
First differences & The changes from one output to the next, on equal input steps. Constant $\to$
    \textbf{linear}. \\
Second differences & The changes in those changes. Constant (when the first are not) $\to$
    \textbf{quadratic}. \\
Common ratio & What each output is multiplied by to reach the next, on equal input steps. Constant
    $\to$ \textbf{exponential}. \\
Extrapolation & Using a model beyond the inputs the data covered. Where a wrong family choice shows
    up --- and where a right one still stops being believable. \\
\end{tabularx}
\end{skillbox}

\begin{fixedskillbox}[Activate Prior Knowledge \& Spiral Review]{mist}
\begin{tabularx}{\linewidth}{@{}X c@{}}
\begin{minipage}[t]{0.55\linewidth}\vspace{0pt}
\textbf{The warm-up is five items, and items 2 and 3 are the plants:}
\begin{enumerate}[leftmargin=1.4em, itemsep=1pt, topsep=2pt]
  \item Last night's item 10 again --- the gym and the bacteria
  \item The adding test: fill a row of first differences (Lesson 2.0)
  \item The multiplying test: fill a row of ratios (Lesson 2.0)
  \item Match the three parents to their shapes (Lesson 2.2)
  \item Which parent has a horizontal asymptote, and what is its equation? (Lesson 2.5)
\end{enumerate}
\end{minipage}
&
\begin{minipage}[t]{0.38\linewidth}\vspace{0pt}
\textbf{How to use the results:} items 2 and 3 are the two tests with the vocabulary withheld --- a
class that can fill both rows in ninety seconds will spend notes box 1 on the \emph{decision} rather
than on arithmetic. Item 5 is the answer to the last row of box 2's comparison table, and item 1
tells you who is still guessing from the picture rather than from the numbers.
\end{minipage}
\end{tabularx}
\end{fixedskillbox}

\begin{skillbox}[Hook]{mist}
\textbf{Two students, one table, and both of them are right about what they checked.} A school video
gets $3$, $6$, $12$, $24$ thousand views on days $0$ through $3$. \quad Student A: \textbf{``The
views go up by $3$, then $6$, then $12$. The jumps keep getting bigger, and a parabola's jumps get
bigger too --- so it's quadratic.''} \ Student B: \textbf{``Every day is double the day before --- so
it's exponential.''} \textbf{Pose it this way:} ``Everything both of them said about these numbers is
true. The jumps really do get bigger, and each value really is double the last. So one of them
checked something that was not \emph{enough}.'' \textbf{Do not resolve it.} 2.6's hook turned on
\emph{where a rule lives}; this one turns on \emph{what counts as evidence}, which is the only
genuinely new idea in the lesson. Settle it on notes box 1, after the three tables.
\end{skillbox}

\boxguard
\begin{skillbox}[Lesson]{mist}
\begin{multicols}{2}
\textbf{1. Two tests, run in order (16 min).} Fill all three tables before saying a single family
name --- the pattern only shows up when all three are on the page at once. Then write the rule, then
the condition (\textbf{equal} input steps; say it out loud, it is the one that gets skipped). Only
then run the tests on the video table and \textbf{resolve the hook}. Close on what Student A actually
checked: ``the jumps get bigger'' rules out \emph{linear} and nothing else.

\textbf{2. The comparison sheet (12 min).} This is the summary of the whole unit and students should
keep it. Fill it fast --- every cell is something they already know from 2.2, 2.4, or 2.5. The two
rows that pay: ``\emph{plus \$30 a month}'' versus ``\emph{up $30\%$ a month},'' and \emph{look for a
turning point first}.

\textbf{3. Using the model, and its window (12 min).} $V(5)=96$ algebraically, then the $50$-thousand
crossing graphically --- \textbf{do both, and name which one you are doing}, because that is
\texttt{AF.1e}. Then day $20$: $3.1$ billion views from a school video. \textbf{The model is not
broken; it is being read outside the days the data covers.}

\textbf{4. Guided practice --- three measurements (10 min).} One of each family, and item 4 is the one
that earns the time: two decreasing data sets, two readings each, and no way to tell them apart.
\end{multicols}
\end{skillbox}

\boxguard
\begin{skillbox}[Explicit Instruction: Choosing a Model]{mist}
\begin{tabularx}{\linewidth}{@{}X|X@{}}
\begin{minipage}[t]{\linewidth}\vspace{0pt}
\textbf{The one-page summary students should be able to reproduce:}
\begin{enumerate}[leftmargin=1.4em, itemsep=1pt, topsep=2pt]
  \item \textbf{Check the input steps first.} Equal, or the tests mean nothing.
  \item \textbf{First differences constant?} $\to$ linear. Stop.
  \item \textbf{Second differences constant?} $\to$ quadratic. Stop.
  \item \textbf{Ratios constant?} $\to$ exponential.
  \item \textbf{From a graph, name a feature:} straight edge $\to$ linear; one turning point $\to$
        quadratic; no turn but a flattening end $\to$ exponential.
  \item \textbf{From a story, listen for the verb:} adds the same \emph{amount} $\to$ linear;
        multiplies by the same \emph{factor or percent} $\to$ exponential; turns around $\to$
        quadratic; ``for the first \dots\ then'' $\to$ piecewise.
  \item \textbf{Then say where it works.} What is the window, and does the prediction pass the
        laugh test?
\end{enumerate}
\textbf{Say this out loud:} ``Never name a family without naming the test.''
\end{minipage}
&
\begin{minipage}[t]{\linewidth}\vspace{0pt}
\textbf{The four errors to name before they happen:}
\begin{itemize}[leftmargin=1.1em, itemsep=2pt, topsep=2pt]
  \item \textbf{``It curves up, so it's exponential.''} The hook, the exit ticket's item 3, Tier E
        item 1, homework item 1(b). \emph{Quadratics curve up too.}
  \item \textbf{Stopping at the first differences.} A student who sees $3, 6, 12$ growing and calls
        it quadratic never checked whether the \emph{second} differences were constant.
  \item \textbf{Running the tests on unequal steps.} Watch for it on the caffeine table (steps of
        $5$) and the braking table (steps of $10$) --- those are fine; a table stepping $1, 2, 5, 10$
        is not.
  \item \textbf{Extrapolating without saying so.} The video at day $20$, the truck at year $8$, the
        shop at month $12$, the tree at year $10$.
\end{itemize}
\end{minipage}
\end{tabularx}
\end{skillbox}

\begin{skillbox}[Active Monitoring]{redbox}
\begin{itemize}[leftmargin=1.2em, itemsep=2pt, topsep=2pt]
  \item \textbf{Notes box 1, table (a):} the ratio row for the rideshare comes out $1.67, 1.4, 1.29$.
        Students want to round those to ``about the same.'' \emph{They are not the same} --- that is
        the point of computing all three.
  \item \textbf{Notes box 1, the hook table:} listen for ``the second differences are $3$ and $6$, so
        it \emph{is} constant-ish.'' Ask: ``Is $3$ equal to $6$?''
  \item \textbf{Notes box 3:} after the day-$20$ number, ask two students ``so is the model wrong?''
        The answer to insist on is \emph{no --- it is being used outside its window}.
  \item \textbf{Activity Tier A item 5:} groups will pick a model. The item does not have a
        model-picking answer --- it has a \emph{measurement} answer. Prompt: ``What would you go
        measure?''
  \item \textbf{Cold-call prompts:} ``Which row was constant?'' ``What feature told you?'' ``Adding
        or multiplying?'' ``Would you bet money on that prediction?''
\end{itemize}
\end{skillbox}

\boxguard
\begin{skillbox}[Group Work \& Differentiation]{redbox}
\begin{multicols}{3}
\textbf{Tier R --- Remediate}
\begin{itemize}[leftmargin=1em, itemsep=1pt, topsep=2pt]
  \item Two duckweed patches: one adds $4$ m$^2$, one triples.
  \item Fill the difference and ratio rows for both.
  \item The week the multiplier passes the adder.
  \item Match each patch to its graph, and name the feature.
\end{itemize}

\columnbreak
\textbf{Tier A --- Approaching Proficiency}
\begin{itemize}[leftmargin=1em, itemsep=1pt, topsep=2pt]
  \item Three graphs, three stories, and a feature for each.
  \item A \$17 fare, read off a graph \emph{and} solved.
  \item $h(10)=-960$ feet --- and why that is not an algebra mistake.
  \item Two points, two models, and the one measurement that settles it.
\end{itemize}

\columnbreak
\textbf{Tier E --- Extension}
\begin{itemize}[leftmargin=1em, itemsep=1pt, topsep=2pt]
  \item Knock down ``it curves, so it's exponential.''
  \item A manager who made \emph{two} errors in one prediction.
  \item The coffee, whose ratios are not constant --- until you subtract the room.
\end{itemize}
\end{multicols}
\end{skillbox}

\boxguard[16]
\begin{skillbox}[Individual Work \& Assessment]{redbox}
\textbf{Exit ticket (3 items, no notes):} a gym counting members over five months --- $64$, $96$,
$144$, $216$, $324$. (1) Fill the three test rows, name the family, name the row that decided it
(\texttt{AF.1c}); (2) predict months $5$ and $6$ ($486$ and $729$) against a building licensed for
$500$, and say how far ahead the model can be trusted (\texttt{AF.1e}); (3) \textbf{the conceptual
check} --- a classmate says the gains speed up, so the model is quadratic.

\textbf{Using the results:} items 1 and 3 are built on the same observation on purpose --- the
2.4/2.5/2.6 pattern used a fourth time. Noticing that the jumps grow is \emph{correct} in item 1 and
the \emph{wrong conclusion} in item 3, so a student who agrees with the classmate because it matches
what they just wrote has told you exactly what they are doing. Sort into three piles: \emph{ran the
ratio test}, \emph{named exponential with no test cited}, and \emph{agreed with the classmate}. The
third pile needs the tests re-taught before the unit test; the second pile needs only the habit.
\textbf{Item 2 is the second read} --- accept any answer that notices $729$ people cannot fit in the
building.
\end{skillbox}

\boxguard[18]
\begin{skillbox}[Reinforcement \& Extension]{goldbox}
\textbf{Homework} (10 items): five fluency items --- three tables to test, four stories with no
numbers, four graphs to classify by feature, the whole comparison sheet from memory, and a full
characteristics read of a fence-and-pen quadratic --- then a four-part read of \textbf{two delivery
trucks depreciated two different ways}, which is the best context in the unit for
\texttt{AF.1c}, \texttt{e}, and \texttt{g} at once.

\textbf{Items 6--9 are the argument for the lesson.} Both trucks cost \$24{,}000. L loses a fixed
\$3{,}000 a year (linear); E loses a fixed $20\%$ a year (exponential). L is worth more in years $1$
through $5$, E is worth more from year $6$ on, and at year $8$ L's model says \textbf{\$0} --- which
no working truck is --- while E's graph approaches \$0 and never reaches it. That is Lesson 2.5's
asymptote arriving in a bank statement.

\textbf{Item 9 is the interpret item of the set.} The buyer's numbers are right and the
\emph{question} is wrong: ``holds its value better'' has no answer until you say \emph{for how long}.
Grade it on whether the answer names a time horizon.

\textbf{Item 10 previews Unit 3:} growth or decay, and the number you multiply by. Two of the four
fall toward an asymptote at $y=0$.

\textbf{Extension (optional):} a tree measured at $3$, $6$, $12$ feet, and two models --- a quadratic
and an exponential --- that fit \textbf{all three points exactly}. They differ by $3$ feet at year $3$
and by $2{,}904$ feet at year $10$. \textbf{Three points can always be fitted by some quadratic}, so
agreeing with the data is not evidence.

\textbf{After this lesson Unit 2 is complete.} The homework's closing box points students at the unit
test and then at \textbf{Unit 3, Exponential and Logarithmic Functions}, where today's constant-ratio
family gets a unit of its own.
\end{skillbox}

% ── Teacher notes — one per component, in packet order ────────────────────────
% Teacher-only prose lives HERE, never in a _key: a note is the one block in a
% key with no counterpart in the blank, so it makes the key longer than the
% blank. See references/conventions.md ("teachernote").
\boxguard[18]
\begin{teachernote}[Warm-Up]
    \textbf{7 minutes, then score it together.} Answers: (1) (a) linear --- the same \$30 is added
    each month; (b) exponential --- the count is multiplied by $2$ each hour; (2) $+4, +4, +4, +4$,
    yes, constant; (3) $3, 3, 3$, yes, constant; (4) $f \to$ (B), $g \to$ (C), $h \to$ (A); (5) only
    $h(x)=2^x$, asymptote $y=0$. \textbf{Items 2 and 3 are the whole lesson with the vocabulary
    withheld}, and they are deliberately easy --- if the room needs more than ninety seconds on them,
    notes box 1 will run long, because box 1 asks for the same two rows three times over. \textbf{Do
    not name ``first differences'' or ``common ratio'' while scoring}; call them ``the change'' and
    ``the ratio'' and let box 1 supply the words. Item 4 is worth watching: a student who matches
    $g(x)=x^2$ to (A) rather than (C) is thinking ``curved'' rather than ``turns around,'' which is
    exactly the error the hook is about to exploit. Item 5's answer is the last row of the comparison
    table in notes box 2, so it will be quoted back within twenty minutes.
\end{teachernote}

\boxguard[18]
\begin{teachernote}[Guided Notes]
    \textbf{Fill all three tables in box 1 before anyone says a family name.} The rule the students
    write is only visible when the three are side by side --- one table has a constant first-difference
    row, one a constant second-difference row, one a constant ratio row, and every \emph{other} row in
    all three tables is not constant. If you let the class classify table (a) as soon as it is filled,
    you lose the pattern. \textbf{The condition on equal input steps is not a footnote} --- table (b)
    steps by $10$ and the caffeine table in the practice steps by $5$, and both are fine; a table
    stepping $1, 2, 5, 10$ is not. Say the difference out loud. \textbf{Resolve the hook only after
    the video table is filled}, and resolve it in these words: the second differences are $3$ then
    $6$, so they are \emph{not} constant, and the ratios are all $2$. Then the sentence that matters
    more than the answer: \textbf{``the jumps get bigger'' rules out linear and nothing else}. Put it
    on the board and leave it there --- it is the exit ticket. Box 2 is the unit's summary sheet;
    tell students to keep it for the test. It fills fast because every cell is 2.2, 2.4, or 2.5. In
    box 3, \textbf{name which method you are using out loud each time} --- ``this one is algebraic,
    this one is graphical'' --- because \texttt{AF.1e} asks for both and students do not notice they
    have done two different things. The day-$20$ number ($3.1$ billion views) is the moment of the
    lesson; ask ``is the model wrong?'' and hold out for \emph{no, it is being used outside its
    window}. If time runs short, assign practice item 3 as reading, but never cut item 4 --- two
    decreasing data sets and two readings each is the cleanest statement of ``the data you have does
    not settle it.''
\end{teachernote}

\boxguard[18]
\begin{teachernote}[Group Activity]
    \textbf{Groups of three, 15 minutes, all groups start at Tier R.} Tier R is one comparison run
    four ways and a fluent group clears it in six minutes; the item that slows them is item 3, where
    Patch A is bigger for two weeks and then loses forever. That is 2.0's adding-versus-multiplying
    with a winner. \textbf{Tier A item 1 is the fastest diagnostic in the lesson} --- if a group
    writes ``(b) is quadratic because it's curved,'' they have matched the shape without naming the
    feature, and the feature is \emph{one turning point}. Graph (c) is curved too. \textbf{Tier A item
    3 is the one worth the argument}: $h(10)=-960$ feet is correct algebra and a meaningless answer,
    because the model stopped applying when the ball hit the ground at $t=4$. \textbf{Tier A item 5 is
    the stretch and most groups will answer the wrong kind of question} --- they will pick a model.
    The item has no model-picking answer; it has a \emph{measurement} answer, and the prompt that
    fixes it is ``what would you go and measure?'' \textbf{Tier E item 3 is the best item in the
    activity and the hardest.} The coffee's ratios are not constant, which looks like a failed test
    --- until you subtract room temperature and find that the \emph{gap} halves exactly. That is a
    vertical translation (2.2) carrying the asymptote up with it (2.5), and it is the honest answer to
    ``why does the ratio test sometimes fail on real data.'' Share it out even if only one group gets
    there.
\end{teachernote}

\boxguard[18]
\begin{teachernote}[Exit Ticket]
    \textbf{5 minutes, notes closed, hand it to me at the door.} Answers: (1) first differences
    $32, 48, 72, 108$; second differences $16, 24, 36$; ratios all $1.5$. \textbf{Exponential},
    decided by the ratio row --- a $50\%$ gain each month. (2) Month $5$: $486$; month $6$: $729$.
    Month $6$ is the first over $500$, and the building cannot hold $729$ people, so the model is
    already predicting something impossible two months past the data. (3) Growing jumps only rule out
    linear; the second differences $16, 24, 36$ are not constant; the ratios are, so the model is
    exponential. \textbf{Items 1 and 3 are built on the same observation on purpose}, the pattern 2.4,
    2.5 and 2.6 all used: noticing the growing jumps is \emph{correct} in item 1 and the \emph{wrong
    conclusion} in item 3. A student who agrees with the classmate because it matches what they wrote
    twelve lines earlier has told you exactly what they are doing. \textbf{Grade item 1 on whether a
    test is cited at all} --- ``exponential'' with no row named is half an answer and should be
    corrected on the spot, because the unit test will ask for the justification. Item 2 needs only
    that the student notices $729$ people will not fit; the arithmetic is the easy part. Do not grade
    for points; use it to decide whether the tests need one more pass before the unit test.
\end{teachernote}

\boxguard[18]
\begin{teachernote}[Homework]
    \textbf{Expect 30--35 minutes.} Item 3 is the fastest to grade and the most diagnostic: a student
    who writes ``exponential'' under graph (ii) has classified by curvature rather than by the turning
    point. Item 4 is the comparison sheet from memory and is worth reminding students about before the
    unit test --- it is the single densest revision page in the unit. \textbf{Items 6--9 are the
    argument for the whole lesson}, and item 8 is where to spend your attention: Truck L's model says
    the truck is worth \textbf{\$0} at year $8$, which is not true of any truck that still runs, while
    Truck E's graph approaches \$0 and never gets there. Students have met that sentence twice --- as
    the coffee gap in 2.5 and as the open circle in 2.6 --- and this is the first time it decides a
    real question. \textbf{Item 9 must be graded on the reasoning, not the verdict}: the buyer's
    numbers are right, and the fault is that ``holds its value better'' has no answer until you name a
    time horizon. Accept any answer that names one. Item 10 is the Unit 3 on-ramp; ``multiplied by
    $1.10$'' and ``multiplied by $0.70$'' are what to look for, and a student who writes ``$10\%$''
    instead of ``$1.10$'' has not yet seen that the multiplier includes the original amount --- fix
    that now, it is the first thing Unit 3 needs. \textbf{The extension is genuinely optional and is
    the sharpest item in the set}: two models fit all three measurements \emph{exactly} and differ by
    almost three thousand feet at year $10$. Students who finish it understand something most adults
    do not, which is that agreeing with the data is not evidence.
\end{teachernote}
\end{document}
